%% sections/03-pipeline.tex

\section{The Style Emergence Pipeline}
\label{sec:pipeline}

\subsection{Promotion Threshold}

A player style is promoted when 3+ innovation log entries cluster thematically across
2+ sessions. The threshold is more demanding than the rubric amendment threshold
(2+ in a dimension): it requires cross-session evidence that the pattern is not
session-specific and cross-instance evidence that it is not a single-event anomaly.

Additionally, the 3 instances must be \textit{substantively distinct} --- not 3
instances of the same mechanic in different scenes, but 3 instances that show the
pattern from different angles or contexts. The pattern's generality is the threshold,
not its frequency.

\subsection{Entry Format}

Each player-style entry has:
\begin{description}
  \item[Definition:] What the behavior is, described in PC-design terms (not player-personality terms).
  \item[Triggers/Signals:] Observable session-log indicators that the style is active in a session.
  \item[Cross-applicability:] Design features that amplify or undermine the style.
  \item[Play implications:] Design guidance for adventure-designers creating content for PCs with this style.
  \item[Confirmed sub-types:] Named variants that have emerged from 3+ cross-campaign instances.
  \item[Cross-campaign validation status:] Which campaigns have produced instances.
\end{description}

\subsection{Validation Protocol}

After initial promotion, a style is monitored across subsequent campaigns for:
(a) instance recurrence (does the pattern continue to fire?); (b) archetype independence
(does it fire in campaigns with different motivational structures?); and (c) sub-type
emergence (do consistent variations accumulate to warrant sub-type naming?).

A style with 3 initial instances and 0 subsequent instances is not withdrawn but is
marked as potentially archetype-specific until cross-campaign evidence accumulates.
No style in the corpus was in this state at the end of the 7-campaign observation period.
